% Part: many-valued-logic
% Chapter: infinite-valued-logics
% Section: introduction

\documentclass[../../../include/open-logic-section]{subfiles}

\begin{document}

\olfileid{mvl}{inf}{int}

\olsection{परिचय}

मॅट्रिक्समधील सत्यतामूल्यांची संख्या ससीम असण्याची गरज नाही.
अनंत सत्यतामूल्यांसाठी एक नैसर्गिक पर्याय म्हणजे $0$ आणि~$1$
यांदरम्यानच्या परिमेय संख्यांचा संच,
$V_\infty = [0,1] \cap \Rat$; म्हणजे
\begin{align*}
    V_\infty & = \Setabs{\frac{n}{m}}{n,m \in \Nat \text{ आणि } 0<m \text{ आणि } n\le m}.
\intertext{हा अनंत सत्यतामूल्यांचा संच विचारात घेताना दोन किंवा
अधिक मूल्ये असलेले पुढील उपसंचही पाहणे अनेकदा उपयोगाचे असते:}
V_m & = \Setabs{\frac{n}{m-1}}{n \in \Nat \text{ आणि } n<m}
\intertext{उदाहरणार्थ, $V_5$ मध्ये समान अंतरावर असलेली $5$ सत्यतामूल्ये आहेत:}
V_5 & = \{0, \frac{1}{4}, \frac{1}{2}, \frac{3}{4}, 1\}.
\end{align*}
या सत्यतामूल्यांच्या संचांवर आधारलेल्या तर्कशास्त्रांत
सहसा फक्त $1$ निर्दिष्ट असते, म्हणजे $V^+ = \{1\}$.
दुसऱ्या शब्दांत, $1$ ला निरपेक्ष सत्याची आणि $0$ ला
निरपेक्ष असत्यतेची भूमिका देतो; पण !!{formula}s ना~$V$
मधील कोणतेही मधले मूल्य मिळू शकते.

याशिवाय $0$ आणि~$1$ यांदरम्यानच्या सर्व \emph{वास्तव}
संख्यांचा संच $V_{[0,1]} = [0,1]$, किंवा $[0,1]$ चे इतर
अनंत उपसंचही विचारात घेता येतात. हा सत्यतामूल्यांचा संच
वापरणाऱ्या तर्कशास्त्रांना अनेकदा \emph{अस्पष्ट} म्हणतात.


\end{document}
